% =====================================================================
%  COMP532 - Assignment 2 Report
%  Deep Reinforcement Learning for LunarLander + Chess Extension
%  Compile with:  pdflatex report.tex (twice for refs)
% =====================================================================
\documentclass[11pt,a4paper]{article}

% --- packages --------------------------------------------------------
\usepackage[utf8]{inputenc}
\usepackage[T1]{fontenc}
\usepackage{lmodern}
\usepackage[a4paper, margin=2.4cm]{geometry}
\usepackage{amsmath, amssymb}
\usepackage{graphicx}
\usepackage{booktabs}
\usepackage{caption}
\usepackage{xcolor}
\usepackage{listings}
\usepackage{fancyhdr}
\usepackage{enumitem}
\usepackage{microtype}
\usepackage[hidelinks, colorlinks=true,
            linkcolor=blue!50!black,
            urlcolor=blue!50!black,
            citecolor=blue!50!black]{hyperref}
\usepackage{titlesec}

% --- look & feel -----------------------------------------------------
\definecolor{primary}{HTML}{1F3864}
\definecolor{accent}{HTML}{2E75B6}
\definecolor{codebg}{HTML}{F4F4F4}
\definecolor{codekw}{HTML}{B11A1A}
\definecolor{codestr}{HTML}{2C7A2C}
\definecolor{codecmt}{HTML}{6F6F6F}

\titleformat{\section}{\color{primary}\Large\bfseries}{\thesection}{0.6em}{}
\titleformat{\subsection}{\color{primary}\large\bfseries}{\thesubsection}{0.6em}{}
\titleformat{\subsubsection}{\color{accent}\normalsize\bfseries}{\thesubsubsection}{0.6em}{}

\pagestyle{fancy}
\fancyhf{}
\fancyhead[L]{\small COMP532 -- Assignment 2}
\fancyhead[R]{\small Page \thepage}
\renewcommand{\headrulewidth}{0.4pt}

\lstset{
  basicstyle=\ttfamily\footnotesize,
  backgroundcolor=\color{codebg},
  keywordstyle=\color{codekw}\bfseries,
  stringstyle=\color{codestr},
  commentstyle=\color{codecmt}\itshape,
  showstringspaces=false,
  breaklines=true,
  frame=none,
  framesep=4pt,
  xleftmargin=8pt,
  language=Python,
  numbers=none,
}

\captionsetup{font=small, labelfont=bf, justification=centering}

\setlist[itemize]{itemsep=2pt, topsep=4pt, parsep=0pt}

% --- title page data -------------------------------------------------
\title{\vspace{-1.5cm}\color{primary}
       \textbf{Deep Reinforcement Learning for LunarLander}\\[2pt]
       \large with a chess-playing extension}
\author{COMP532 -- Assignment 2 Report}
\date{}

\begin{document}

% =====================================================================
\maketitle\thispagestyle{fancy}
\vspace{-1cm}

\begin{center}
\textit{A Dueling Double-DQN agent trained from scratch in PyTorch, plus a
random / heuristic / DRL / LLM ablation on chess.}
\end{center}

\section*{Group members and contributions}
\noindent\textit{Please replace the table below with your group's
details before submission. Identical marks are awarded to all members
of a group, and every student must submit the report individually
via Canvas.}

\vspace{4pt}
\begin{center}
\small
\begin{tabular}{p{4cm}p{4cm}p{6cm}}
\toprule
\textbf{Name} & \textbf{Email / Student ID} & \textbf{Contribution} \\
\midrule
$\langle$Student 1$\rangle$ & $\langle$email / id$\rangle$ & DDDQN agent code, training loop, LunarLander runs.\\
$\langle$Student 2$\rangle$ & $\langle$email / id$\rangle$ & Replay buffer, plotting utilities, evaluation harness.\\
$\langle$Student 3$\rangle$ & $\langle$email / id$\rangle$ & Chess DRL value-net, board encoding, training loop.\\
$\langle$Student 4$\rangle$ & $\langle$email / id$\rangle$ & LLM agent, prompt design, Stockfish wrapper, UCI shim.\\
$\langle$Student 5$\rangle$ & $\langle$email / id$\rangle$ & Report writing, figures, tournament harness, references.\\
\bottomrule
\end{tabular}
\end{center}

\vspace{4pt}
\noindent\textbf{Reproducibility.} A single command reproduces every number in this
report (training is $\approx 30$ minutes on CPU):
\begin{lstlisting}[language=bash]
python main.py --episodes 1000 --seed 42                    # Problem 1
python -m chess_ext.train_drl --n-games 300 --seed 42       # extension
python -m chess_ext.run_tournament --n-games 10 --seed 123  # ablation
\end{lstlisting}

\vspace{4pt}
\noindent\textbf{Headline results.}
\begin{center}\small
\begin{tabular}{lr}
\toprule
\textbf{Metric} & \textbf{Value} \\
\midrule
LunarLander solved at episode (mean$_{100}\!\geq\!200$) & \textbf{462} \\
LunarLander best 100-ep mean reward & 252.87 \\
LunarLander greedy evaluation (30 unseen seeds) & $228.80 \pm 69.41$ \\
Chess DRL best combined win-rate vs (Random, Heuristic) & \textbf{75.0\%} \\
Tournament: DRL win-rate over 40 games & 42.5\% (12-18-10) \\
Tournament: LLM-stub win-rate over 40 games & 55.0\% (15-11-14) \\
\bottomrule
\end{tabular}
\end{center}

% =====================================================================
\newpage
\section{Problem 1: Deep RL agent for LunarLander}

\subsection{Environment}
LunarLander is a Box2D physics task in which an under-actuated lander
must softly touch down on a randomly generated landing pad
\cite{gymnasium}. The state $s\in\mathbb{R}^{8}$ contains the lander's
$(x, y)$ position, velocity $(v_x, v_y)$, orientation $\theta$, angular
velocity $\omega$, and two booleans for left/right leg contact. The
action space is discrete with four primitives: do nothing, fire left
engine, fire main engine, fire right engine. Rewards combine dense
shaping (proximity, velocity, tilt, leg contact, fuel use) with a
terminal $+100$ for a soft landing on the pad and $-100$ for a crash.
The task is conventionally considered \emph{solved} when the average
return over 100 consecutive episodes exceeds 200.

We use \texttt{LunarLander-v3} from Gymnasium (\textgreater{}=1.0). The
underlying environment is identical to the older \texttt{LunarLander-v2}
shipped by OpenAI Gym; only the version tag changed. Our training
harness tries v3 first and falls back to v2.

\subsection{Choice of algorithm}
\label{sec:alg}
We implement a \textbf{Dueling Double DQN (D3QN)}. This is a controlled
upgrade of the vanilla DQN suggested in the assignment starter and
combines two well-established improvements:

\begin{itemize}
  \item \textbf{Double DQN} \cite{vanhasselt2016} decouples action
    selection from action evaluation in the bootstrap target. Standard
    DQN uses $\max_a Q_{\text{tgt}}(s', a)$, which over-estimates $Q$
    because the same network both selects and evaluates the maximising
    action. Double DQN selects with the online network and evaluates
    with the frozen target,
    \[
        y \;=\; r + \gamma\,(1-\text{done}) \cdot
                Q_{\text{tgt}}\!\bigl(s',\, \arg\max_a Q_{\text{online}}(s', a)\bigr),
    \]
    removing most of the maximisation bias.
  \item \textbf{Dueling architecture} \cite{wang2016} factorises
    $Q(s, a) = V(s) + (A(s, a) - \frac{1}{|A|}\sum_{a'} A(s, a'))$. In
    LunarLander, many states are nearly action-indifferent (e.g.\
    coasting toward the pad with the engine off), and explicitly
    learning a state-value channel reduces the variance of the policy
    improvement step.
\end{itemize}

Both modifications are zero-overhead at inference, are invariant to the
standard DQN training loop, and consistently improve sample efficiency
and stability \cite{hessel2018}. We deliberately did not use a
high-level RL library (e.g.\ Stable-Baselines3) so that all
algorithmic details remain transparent and assessable. Policy-gradient
or actor-critic methods (REINFORCE, A2C, PPO) would have been viable
alternatives, but value-based methods are conventional for small
discrete-action tasks and integrate naturally with experience replay.

\subsection{What we changed relative to the starter notebook}
The provided starter implements vanilla DQN with a $64\!\to\!64$ MLP
trunk and ``fixed Q-targets'' (periodic hard copy). Our submission is a
strict superset; the upgrades and their motivations are summarised in
Table~\ref{tab:starter-vs-ours}.

\begin{table}[h]
\centering\small
\caption{Differences between the starter notebook and our submission.}
\label{tab:starter-vs-ours}
\begin{tabular}{lll}
\toprule
\textbf{Aspect} & \textbf{Starter notebook} & \textbf{Our submission} \\
\midrule
Algorithm & DQN with fixed Q-targets & Dueling Double DQN \\
Trunk width & $64\!\to\!64$ & $128\!\to\!128$ \\
Heads & single output of size $|\mathcal{A}|$ & $V(s)$ + $A(s, a)$ \\
Target update & periodic hard copy (every $C$ steps) & Polyak soft, $\tau\!=\!10^{-3}$ \\
Replay & deque of tuples & pre-allocated NumPy ring buffer \\
Loss & MSE & Smooth-L1 (Huber) + grad-clip 10 \\
Bootstrap & $\max_a Q_{\text{tgt}}$ & Double-Q target (Sec.~\ref{sec:alg}) \\
\bottomrule
\end{tabular}
\end{table}

\subsection{Network architecture}
A small MLP is adequate for the 8-D LunarLander state. The shared
feature trunk has two hidden layers of 128 ReLU units. Two heads then
branch off:
\begin{itemize}
  \item \textbf{Value head $V(s)$}: $\text{Linear}(128) \to \text{ReLU} \to \text{Linear}(1)$.
  \item \textbf{Advantage head $A(s, a)$}: $\text{Linear}(128) \to \text{ReLU} \to \text{Linear}(4)$.
\end{itemize}
The Q-values are recombined with the standard mean-subtracted
aggregator $Q(s, a) = V(s) + A(s, a) - \tfrac{1}{4}\sum_{a'} A(s, a')$.
Mean-subtraction (rather than max) is preferred in practice because it
is differentiable everywhere and yields a more stable gradient signal.
Linear layers are He-initialised because of the ReLU activations.

\subsection{Replay, target network, and exploration}
\begin{itemize}
  \item \textbf{Experience replay} uses pre-allocated NumPy ring
    buffers $(s, a, r, s', \text{done})$ with capacity 100{,}000.
    Pre-allocation eliminates per-step allocation overhead, giving
    $\approx 5\times$ the throughput of a Python deque-of-tuples
    implementation.
  \item \textbf{Target network} is updated by Polyak (soft) averaging
    with $\tau = 10^{-3}$ every learner step. This produces a smoothly
    moving target without the abrupt distribution shifts of periodic
    hard copies, while still lagging the online network by a factor of
    ${\approx}1000$ steps.
  \item \textbf{Exploration} uses linearly-decaying $\varepsilon$-greedy:
    $\varepsilon$ starts at 1.0 and is multiplied by 0.995 after every
    episode until it reaches a floor of 0.01. Per-episode decay (rather
    than per-step) is convenient because LunarLander episode lengths
    vary by an order of magnitude during training.
  \item \textbf{Learning step} is taken every 4 environment steps once
    the buffer holds at least 1{,}000 transitions. The Huber loss is
    used because the $\pm100$ terminal reward creates occasional
    outliers that would distort an MSE objective. Gradients are clipped
    to L2 norm 10 to prevent exploding updates early in training.
\end{itemize}

\subsection{Hyper-parameters}
\begin{table}[h]
\centering\small
\caption{All hyper-parameters used. Defaults follow common DQN
practice; only $\varepsilon$-decay and replay capacity were tuned on a
100-episode pilot.}
\label{tab:hp}
\begin{tabular}{p{4cm}p{2.5cm}p{8cm}}
\toprule
\textbf{Hyper-parameter} & \textbf{Value} & \textbf{Justification} \\
\midrule
Algorithm & Dueling DDQN & Combines Double-DQN bias correction with Dueling V/A decomposition. \\
Discount factor $\gamma$ & 0.99 & Standard for episodic continuous-control tasks. \\
Learning rate & $5\!\times\!10^{-4}$ & Adam default region for small MLP; empirically stable on LunarLander. \\
Batch size & 64 & Common DQN choice; balances variance and throughput. \\
Replay capacity & 100{,}000 & Holds $\approx500$ episodes of experience to break sample correlations. \\
Warm-up steps & 1{,}000 & Avoids learning from a near-empty buffer. \\
Target update & Polyak $\tau\!=\!10^{-3}$ & Smoother than periodic hard copy; matches DDPG/Rainbow practice. \\
Update frequency & every 4 env steps & Decouples acting from learning. \\
Network width & $128\times 2$ + heads & Sufficient for an 8-D state $\times$ 4-action problem. \\
Loss & Smooth-L1 (Huber) & Robust to occasional large reward outliers ($\pm100$). \\
Gradient clip & L2 norm 10 & Prevents exploding gradients early. \\
$\varepsilon$ schedule & $1.0\!\to\!0.01$, decay 0.995/ep & Smooth annealing reaches floor near ep 920. \\
Random seed & 42 & Set on torch, numpy and the env reset. \\
\bottomrule
\end{tabular}
\end{table}

\subsection{Implementation outline}
The codebase is split into small, single-purpose modules to make the
algorithm easy to read and audit:
\begin{lstlisting}
lunar_lander_drl/
  src/
    network.py        # DuelingQNetwork
    replay_buffer.py  # ReplayBuffer (NumPy ring buffer)
    agent.py          # DuelingDoubleDQNAgent + AgentConfig
    train.py          # training loop with rolling-mean tracking
    evaluate.py       # greedy roll-outs and GIF recording
    plotting.py       # reward and loss figures
  main.py             # CLI entry point
  chess_ext/          # Optional extension (Section 3)
\end{lstlisting}

\subsection{Results}
The agent was trained for 1000 episodes ($\approx 19.6$ minutes on a
single CPU). The best 100-episode rolling mean reached \textbf{252.87}
and the task was \textbf{solved at episode 462} (rolling mean reaching
the threshold of 200). The final 100-episode mean reward at the end of
training was 246.40.

\begin{figure}[h]
  \centering
  \includegraphics[width=0.85\textwidth]{plots/rewards.png}
  \caption{Episode reward (light) and 100-episode rolling mean (solid)
  over training. The horizontal dashed line marks the conventional
  ``solved'' threshold of 200.}
  \label{fig:rewards}
\end{figure}

\begin{figure}[h]
  \centering
  \includegraphics[width=0.85\textwidth]{plots/loss.png}
  \caption{Per-update Huber loss vs.\ gradient step (log $y$-axis), with
  EMA smoothing for readability. The high-variance early phase
  ($\approx0$--20k steps) coincides with the agent first encountering
  successful landings, which inject high-magnitude returns into the
  buffer.}
  \label{fig:loss}
\end{figure}

\begin{figure}[h]
  \centering
  \includegraphics[width=0.55\textwidth]{plots/demo_frame.png}
  \caption{A representative frame from a greedy roll-out of the trained
  agent (also saved as \texttt{videos/agent\_demo.gif}). The lander has
  touched down between the two flags.}
  \label{fig:demo}
\end{figure}

\subsection{Discussion}
Figure~\ref{fig:rewards} shows the three phases that are characteristic
of DQN on dense-reward tasks. (i) For roughly the first 100 episodes
the agent's policy is essentially random because $\varepsilon$ is still
high and the replay buffer is filling; episodic returns hover around
$-150$ because the lander typically crashes or fires fuel uselessly.
(ii) Between episodes ${\approx}100$--350, returns climb steeply: the
buffer is now diverse enough that gradient steps are informative, and
$\varepsilon$ has decayed enough that exploitation dominates.
(iii) After the rolling mean crosses zero, progress slows because
further improvement requires fine motor control to land between the
flags rather than crash near them. Convergence to the solved threshold
typically occurs by episode 400--500; we observed crossing at episode
462.

The loss curve in Figure~\ref{fig:loss} is non-monotonic, which is
expected and not a sign of divergence: as the policy improves, the
distribution of bootstrap targets shifts and the network is effectively
chasing a moving objective. The Polyak-averaged target net and gradient
clipping keep this distribution shift bounded. The Huber loss decays
roughly two orders of magnitude over training but never approaches
zero, because there remains intrinsic stochasticity in the environment
(initial conditions, contact dynamics).

Greedy evaluation ($\varepsilon = 0$) over 30 unseen seeds yields a
mean return of \textbf{228.80} (std 69.41) with 70\% of roll-outs
exceeding the 200 ``solved'' threshold, confirming that the policy
generalises rather than overfitting to specific replayed transitions.

% =====================================================================
\newpage
\section{Problem 2: Exploration vs.\ exploitation in deep RL}

\subsection{Definitions and intuition}
In a reinforcement-learning problem, an agent interacts with an unknown
MDP and must choose actions that maximise expected discounted return.
\textbf{Exploitation} means selecting the action currently estimated to
be best -- i.e.\ $\arg\max_a Q(s, a)$ for the agent's value function.
\textbf{Exploration} means deliberately taking actions whose value is
uncertain in order to gather information that may revise those
estimates. The two are in fundamental tension: a purely exploiting
agent gets stuck in whatever locally optimal behaviour it stumbles into
first, while a purely exploring agent never converts its information
into reward.

The trade-off is usually formalised through the concept of
\emph{regret}: the cumulative expected difference between the optimal
policy's return and the agent's return,
$R(T) = \sum_{t=1}^{T}\bigl(V^{*}(s_t) - V^{\pi_t}(s_t)\bigr)$.
Good exploration strategies have provably sub-linear regret growth,
meaning the per-step performance gap shrinks toward zero as experience
accumulates \cite{sutton2018}.

\subsection{Why the trade-off is unavoidable in deep RL}
Three structural reasons make the exploration--exploitation trade-off
intrinsic to deep RL:
\begin{itemize}
  \item \textbf{Bootstrapping}: deep value-based methods learn from
    their own predictions (via the Bellman backup). If the agent never
    visits a state-action pair, its $Q$-estimate there remains
    arbitrary and the policy can be globally pessimistic about a
    profitable region.
  \item \textbf{Function approximation}: a neural $Q$-network
    generalises across states. Optimistic exploration of one part of
    the space may distort estimates elsewhere through shared
    parameters, which can hurt as well as help -- so exploration must
    be calibrated, not maximal.
  \item \textbf{Sparse or deceptive reward}: when reward is rare or
    shaped misleadingly (e.g.\ early local maxima), naive exploitation
    converges quickly to a sub-optimal policy. Exploration is what lets
    the agent escape such basins of attraction.
\end{itemize}

\subsection{Common strategies}
\subsubsection{$\varepsilon$-greedy}
With probability $\varepsilon$ the agent samples a uniformly random
action; otherwise it picks $\arg\max_a Q(s, a)$. $\varepsilon$ is
typically annealed from 1.0 to a small floor (e.g.\ 0.01 or 0.05) so
that the agent transitions smoothly from exploration to exploitation.
This is the strategy used in our LunarLander implementation. It is
simple, has only one hyper-parameter, and works well when reward is
reasonably dense, but it is \emph{undirected}: every non-greedy action
is equally likely, regardless of how informative it would be.

\subsubsection{Boltzmann / softmax}
Actions are sampled from a softmax over Q-values:
$\pi(a \mid s) \propto \exp\bigl(Q(s, a) / \tau\bigr)$. The temperature
$\tau$ controls the explore--exploit balance and can be annealed
similarly to $\varepsilon$. Softmax exploration weights actions by
their estimated value, so it explores promising-looking suboptimal
actions more than clearly bad ones. The downside is sensitivity to the
absolute scale of $Q$-values, which is why it is less common in deep
value-based methods.

\subsubsection{Upper Confidence Bound (UCB)}
Bandits and tabular MDPs admit principled exploration via confidence
bounds:
$a_t = \arg\max_a \bigl[Q(s, a) + c\sqrt{\ln t / N(s, a)}\bigr]$.
The bonus term shrinks with the visit count $N(s, a)$, so
under-visited actions are tried more often. UCB has provable regret
guarantees but requires explicit count statistics, which are
intractable in continuous or high-dimensional state spaces. Deep RL
approximates this with pseudo-counts derived from density models or
learned hash functions \cite{burda2019rnd}.

\subsubsection{Entropy regularisation}
In policy-gradient methods (e.g.\ A2C, PPO, SAC) exploration is
encouraged by adding a bonus proportional to the policy entropy
$\mathcal{H}\bigl[\pi(\cdot \mid s)\bigr]$ to the objective. This pushes
the policy to remain stochastic until the value evidence forces it to
commit. SAC pushes this further by formulating the entire problem as
maximum-entropy RL. Entropy methods are particularly effective in
continuous action spaces where $\varepsilon$-greedy is ill-defined.

\subsubsection{Noisy networks and parameter-space noise}
Instead of perturbing the chosen action, NoisyNets \cite{fortunato2018}
inject learnable Gaussian noise into the parameters of the final
layers. The noise variance is itself a learned parameter, so the agent
automatically reduces exploration in well-understood states and
increases it elsewhere. This is the exploration mechanism used in
Rainbow DQN.

\subsubsection{Intrinsic motivation and curiosity}
When extrinsic reward is very sparse (e.g.\ Montezuma's Revenge) the
methods above fail because random actions almost never produce useful
feedback. Curiosity-driven approaches add an intrinsic reward
proportional to a prediction error -- typically a forward model's error
on the next state (ICM \cite{pathak2017}, RND \cite{burda2019rnd}) --
so the agent is rewarded for visiting genuinely novel states.

\subsection{Practical challenges}
\begin{itemize}
  \item \textbf{Schedule sensitivity}: the rate at which $\varepsilon$
    (or entropy weight) is annealed is often more important than the
    algorithm itself. Decay too fast and the agent commits before
    sufficient data has been collected; decay too slowly and learning
    is wasted on random behaviour after a good policy is available.
  \item \textbf{Off-policy stability}: aggressive exploration produces
    transitions whose action distribution differs strongly from the
    current greedy policy. With function approximation, this can
    interact badly with bootstrapping and cause divergence (the
    ``deadly triad''). Replay-buffer mixing and target networks
    mitigate but do not eliminate the issue.
  \item \textbf{Reward shaping interactions}: shaped rewards (such as
    LunarLander's proximity term) can mislead exploration by making
    locally rewarding behaviours look optimal. Greedy descent toward
    the centre of the screen earns shaping reward but does not solve
    the task on its own.
  \item \textbf{Reproducibility}: because exploration is stochastic and
    bootstraps interact with function approximation, small differences
    in initial random actions can produce large differences in final
    performance. Reporting variance over multiple seeds is therefore
    standard practice.
\end{itemize}

\subsection{Application to LunarLander}
LunarLander has dense, well-shaped rewards and a low-dimensional state,
so simple $\varepsilon$-greedy exploration is sufficient and was chosen
here. Three properties of our setup were tuned with the trade-off
explicitly in mind:
\begin{itemize}
  \item \textbf{Initial $\varepsilon = 1.0$} for the first $\approx 70$
    episodes guarantees broad state coverage before any greedy
    exploitation, which is important because the early random policy
    is otherwise heavily biased toward firing fuel uselessly.
  \item \textbf{Multiplicative decay 0.995 per episode} reaches the
    floor ($\varepsilon = 0.01$) at episode $\approx 920$. This allows
    the agent to keep $\approx 1\%$ exploration well past the point
    where the rolling mean exceeds 200, which prevents premature
    commitment when later episodes occasionally start in unusual
    configurations.
  \item \textbf{Floor of 0.01} (rather than 0) keeps a small residual
    exploration even after ``solving'', so that any non-stationarity
    introduced by the moving target network does not let the agent
    become deterministically trapped on a sub-optimal landing approach.
\end{itemize}
The reward curve in Figure~\ref{fig:rewards} is consistent with this
design: returns climb steadily once $\varepsilon$ falls below
$\approx 0.4$ (around episode 200), and the loss curve in
Figure~\ref{fig:loss} shows the matching transition from high-variance
updates to a slowly decaying steady state. We did experiment briefly
with softmax exploration on a 100-episode pilot; results were
comparable on this dense-reward task and we therefore retained the
simpler $\varepsilon$-greedy baseline. For sparser tasks (e.g.\
continuous-control suites or Atari), entropy regularisation or
NoisyNets would be more appropriate.

% =====================================================================
\newpage
\section{Optional extension: chess-playing agents}

\subsection{Scope and methodological choices}

The brief offers an optional chess extension, suggesting either a
deep-RL or LLM-based agent. We implemented \textbf{both} so we can
compare them on a level playing field, and we evaluate them in a
five-agent round-robin tournament against the required random
baseline plus a heuristic agent and a weak Stockfish reference. All
agents share a common \texttt{BaseAgent} interface and a tournament
harness so the comparison is reproducible from one command.

\paragraph{Why not OpenAI Gym?} There is no canonical Gymnasium
environment for chess; the older \texttt{gym-chess} package on PyPI is
unmaintained, uses a clunky one-hot action encoding, and is not used in
the published chess-RL literature. AlphaZero \cite{silver2018alphazero},
Leela Chess Zero, ChessGPT \cite{feng2023chessgpt} and the recent
LLM-Chess benchmark \cite{kolasani2025llmchess} all wrap
\texttt{python-chess} \cite{pythonchess} directly. We follow the same
convention: \texttt{python-chess} provides a complete legal-move
generator, FEN/PGN I/O, and a UCI engine wrapper, which is exactly
what the project needs.

\paragraph{Why not chess.com?} The brief asks whether the agent could
play \texttt{chess.com}. Honestly: not via a supported interface.
\texttt{chess.com} provides no public bot/play API. The two ways a
program could play there in practice are (i) authenticated
\emph{browser automation} against the chess.com web UI, and (ii)
reverse-engineering the private REST endpoints used by their own
front-end. Both violate the \texttt{chess.com} Terms of Service and are
account-bannable. The intentional, supported route to ``let an engine
play online against humans'' is \textbf{lichess.org}, which exposes a
first-class \emph{Bot API} (OAuth2 scope \texttt{bot:play}) explicitly
for engine authors. The official bridge software,
\texttt{lichess-bot} \cite{lichessbot}, communicates with engines over
the standard \emph{Universal Chess Interface (UCI)}. We therefore
implemented our agents behind a UCI shim
(\texttt{python -m chess\_ext.uci}) so that any of them can be
deployed on lichess in a few configuration lines (see
Section~\ref{sec:online}).

\subsection{Research summary: DRL vs.\ LLMs for chess}

The DRL-vs-LLM comparison for chess has been studied actively since
GPT-3.5 was first observed playing at amateur club level
\cite{carlini2023chessllm}. Three findings from the literature inform
our design:

\begin{enumerate}
  \item \textbf{Pure DRL is the only paradigm that has reached
    super-human chess.} AlphaZero learned chess from self-play with
    MCTS-guided value/policy networks, using $\approx5{,}000$
    first-generation TPUs over nine hours
    \cite{silver2018alphazero}. No LLM-based system has come close.
    Smaller-scale DRL experiments (e.g.\ \cite{zjeffer2022chesszero})
    typically beat random opponents but plateau well below club level
    on practical compute budgets.
  \item \textbf{LLMs play chess by linguistic pattern-matching, not
    search.} GPT-3.5-turbo-instruct achieves $\approx1800$ Elo on
    lichess by predicting the next move in the way a strong human
    annotator would write it down \cite{carlini2023chessllm}. Modern
    reasoning models (o3, o4-mini) reach Elo $\approx500$--$750$
    against Komodo Dragon at low skill settings, with reasoning effort
    directly translating to playing strength
    \cite{kolasani2025llmchess}. Without chess-specific fine-tuning,
    LLMs frequently hallucinate illegal moves and collapse in the
    endgame because the autoregressive context becomes increasingly
    unreliable.
  \item \textbf{Hybrid approaches dominate.} ChessGPT
    \cite{feng2023chessgpt} fine-tunes a transformer on chess corpora
    and outperforms generalist GPT-4 on binary move-selection tasks.
    Recent work distills value-function feedback into LLM training
    with verifiable rewards, but all such systems plateau ``far below
    expert levels'' \cite{strategicrl2025}.
\end{enumerate}

\paragraph{Recommendation.} Both paradigms are valuable for different
reasons: DRL is the only path to strong play on a finite compute
budget; LLMs offer interpretable, conversational behaviour. We
therefore implement both, evaluate them against the required random
baseline, and let the tournament tell the story. We are explicit that
our DRL agent is \emph{not} AlphaZero -- we trained it for $\approx10$
minutes on CPU; the goal of the extension is to demonstrate the
\emph{methodology} rather than to chase Elo.

\subsection{Five agents under test}
\label{sec:five-agents}

\begin{itemize}
  \item \textbf{RandomAgent} (required baseline): uniform random over
    legal moves.
  \item \textbf{HeuristicAgent}: classical 1-ply greedy with a
    material-plus-piece-square-table evaluation (values 100/320/330/500/900
    for P/N/B/R/Q; piece-square tables from the Chess Programming
    Wiki). For each legal move, push, evaluate, pop; play the move
    that maximises evaluation from the side-to-move's perspective.
    Ties broken at random.
  \item \textbf{LLMAgent}: a configurable agent with two backends
    sharing the same prompt-and-parse pipeline (Section~\ref{sec:llm}).
    The \texttt{stub} backend runs offline and reproduces the
    documented LLM failure modes (hallucinated moves, late-game
    collapse, endgame blindness). The \texttt{openai}/\texttt{anthropic}
    backends issue real API calls.
  \item \textbf{DRLAgent}: a small CNN value network trained by
    temporal-difference learning against a weak Stockfish
    (Section~\ref{sec:drl}).
  \item \textbf{StockfishAgent}: a wrapper around Stockfish 16
    \cite{stockfish}, configured at \emph{Skill Level 0, depth 1} so
    it is weak enough to be trained against in tens of minutes. Used
    as both the DRL training opponent and a tournament reference
    point.
\end{itemize}

\subsection{LLM agent design}
\label{sec:llm}

\paragraph{Prompt and parser.} The agent serialises the position as
FEN, lists all legal moves in SAN, and asks the model to reply with one
SAN move. The system prompt enforces the format:
\begin{lstlisting}[language=]
You are a chess engine. You will be given the current board position in
FEN notation along with the list of legal moves in standard algebraic
notation (SAN). Reply with EXACTLY ONE legal move in SAN. Do not
include any other text, commentary, punctuation, or move number.
Choose the move that is best for the side to move.
\end{lstlisting}
The parser strips common preambles (``My move is'', ``I play'', ``\,Move:''),
takes the first whitespace token, and tries SAN; failing that, it
falls back to UCI. If no parse succeeds, the agent re-queries with
the explicit message that the previous response was illegal, up to
\texttt{max\_retries=2} times. On final failure it falls back to
\texttt{HeuristicAgent}, so the pipeline never forfeits a game on a
bad reply.

\paragraph{Why a stub backend.} A real LLM call cannot be issued from
the sandbox in which the experiments were run. Rather than fake a
constant ``e4'' response that would mis-represent the agent, we built
an offline simulator that reproduces the LLM behaviour reported in
\cite{kolasani2025llmchess, carlini2023chessllm}: opening-book play
in the first $\approx5$ plies, a capture-aware heuristic in the
middle-game with growing noise to mirror late-game collapse, and a
configurable rate of \emph{illegal-move hallucinations}
(\texttt{illegal\_rate=0.05}) so that the legality-repair loop is
genuinely exercised. The same code path supports a real model -- the
user passes \texttt{--backend openai --model gpt-4o-mini} (and an API
key) to issue real calls. This design choice is documented honestly
in the report rather than disguised; the alternative was either
fabricating numbers or omitting the LLM agent entirely.

\paragraph{Reproducibility-affecting limitations of the LLM path.}
\begin{itemize}
  \item Closed-model API responses are non-deterministic at any
    \texttt{temperature\,>\,0}. With \texttt{temperature\,=\,0.2} the
    same query can still yield different moves on different days.
  \item Provider-side model versions are silently updated (e.g.\
    \texttt{gpt-4o-mini} pinned to a date), so re-running months later
    will produce different numbers.
  \item Token-level rate limiting and per-call cost mean that running
    even 100 games can exceed an undergraduate-project budget. Our
    stub costs \pounds 0.
\end{itemize}

\subsection{DRL agent design}
\label{sec:drl}

\paragraph{Architecture.} The DRL agent uses a small CNN
\emph{value network} $V_\theta(s) \in [-1, +1]$ predicting the
expected outcome of a position from the side-to-move's perspective.
Move selection at play time is one-ply look-ahead:
\[
   a^{*} \;=\; \arg\min_{a\in\mathcal{A}_{\text{legal}}(s)}\;
                V_\theta\bigl(\text{push}(s, a)\bigr).
\]
The minimum is taken because after my move the opponent is to move,
so a higher $V_\theta$ from \emph{their} perspective means a worse
outcome for \emph{me}. Two design consequences:
\begin{itemize}
  \item Legal moves are guaranteed by construction -- we only ever
    score positions in the move-generator's output, so the agent
    cannot hallucinate illegal moves.
  \item No fixed move vocabulary is needed (chess has $\approx 4{,}672$
    distinct \emph{(from-sq, to-sq, promotion)} actions in AlphaZero's
    encoding); we side-step this by scoring child states directly.
\end{itemize}

\paragraph{Encoding.} Boards are encoded as $18\!\times\!8\!\times\!8$
tensors: 12 piece-channel planes (white P/N/B/R/Q/K then black), 4
castling-rights planes (filled or empty), 1 side-to-move plane, and 1
en-passant-target plane.

\paragraph{Network.} Three convolutional blocks ($18\to 32\to 32\to 32$
filters, $3\times3$ kernels, ReLU) followed by global average pooling,
a $32\to 64$ ReLU layer, and a $64\to 1$ tanh output. About 25{,}000
parameters; trains on CPU in tens of minutes.

\paragraph{Training.} Self-play against a weak Stockfish (Skill 0,
depth 1) is the simplest tractable signal in the time budget. Each
game's terminal outcome ($+1$/$-1$/$0$ from each side's perspective)
is propagated to every position visited, blended with a heuristic
shaping target $\tanh(\text{eval}/600)$ at weight $\alpha$ that
linearly decays from 0.5 to 0.1 over training. We use Adam at
$\text{lr}=10^{-3}$, batch size 256, 8 gradient steps per game, MSE
loss, gradient clip 1.0, $\varepsilon$-greedy exploration at the agent
($\varepsilon: 0.30\!\to\!0.05$), and a 50{,}000-position replay
buffer. The best snapshot by combined win-rate against
RandomAgent and HeuristicAgent (evaluated every 20 games over 6 games
each) is kept.

\subsection{Chess training results}

The DRL agent was trained for 300 self-play games ($\approx9.8$ minutes
on CPU). The best snapshot achieves a \textbf{combined win-rate of 75\%}
against the two baselines (100\% vs Random, 50\% vs Heuristic).
Figure~\ref{fig:chess-train} shows the win-rate curves; Figure~\ref{fig:chess-loss}
shows the value-network loss.

\begin{figure}[h]
  \centering
  \includegraphics[width=0.85\textwidth]{chess_ext/plots/training.png}
  \caption{DRL chess agent: evaluation win-rate against the Random and
  Heuristic baselines, evaluated every 20 self-play training games.
  Combined = mean of the two. The dashed line is the coin-flip baseline.}
  \label{fig:chess-train}
\end{figure}

\begin{figure}[h]
  \centering
  \includegraphics[width=0.85\textwidth]{chess_ext/plots/loss.png}
  \caption{Value-network MSE loss vs.\ gradient step (linear $y$-axis).
  The slow upward drift after step ${\approx}1{,}500$ reflects the
  value-target distribution shift as the policy improves and visits
  more sharply-decided positions.}
  \label{fig:chess-loss}
\end{figure}

\subsection{Tournament: full ablation}
\label{sec:tournament}

We ran a complete round-robin: every pair of agents plays $N=10$ games
with alternating colours, totalling $\binom{5}{2} \times 10 = 100$
games. Per-pair results are in Table~\ref{tab:tournament}; the
heat-map view is in Figure~\ref{fig:matrix}.

\begin{table}[h]
\centering\small
\caption{Round-robin tournament: each cell is wins--losses--draws
\emph{from the row player's perspective} over 10 games (5 with each
colour). The right-most column is the row player's overall win-rate
(wins\,+\,0.5\,draws) summed across all 40 of their games.}
\label{tab:tournament}
\begin{tabular}{lcccccc}
\toprule
 & \textbf{Random} & \textbf{LLM-stub} & \textbf{Heuristic} & \textbf{DRL} & \textbf{Stockfish-0} & \textbf{Total} \\
\midrule
\textbf{Random}        & --        & 0--9--1 (5\%)  & 0--9--1 (5\%)  & 0--10--0 (0\%) & 0--10--0 (0\%) & \textbf{2.5\%} \\
\textbf{LLM-stub}      & 9--0--1 (95\%) & --        & 1--1--8 (50\%) & 5--0--5 (75\%) & 0--10--0 (0\%) & \textbf{55.0\%} \\
\textbf{Heuristic}     & 9--0--1 (95\%) & 1--1--8 (50\%) & --       & 3--2--5 (55\%) & 0--9--1 (5\%)  & \textbf{51.2\%} \\
\textbf{DRL}           & 10--0--0 (100\%) & 0--5--5 (25\%) & 2--3--5 (45\%) & --   & 0--10--0 (0\%) & \textbf{42.5\%} \\
\textbf{Stockfish-0}   & 10--0--0 (100\%) & 10--0--0 (100\%) & 9--0--1 (95\%) & 10--0--0 (100\%) & --  & \textbf{98.8\%} \\
\bottomrule
\end{tabular}
\end{table}

\begin{figure}[h]
  \centering
  \includegraphics[width=0.85\textwidth]{chess_ext/plots/tournament_matrix.png}
  \caption{Head-to-head win-rate as a heat-map: each cell shows
  wins--losses--draws (and per-cent score) for the row player vs the
  column player over 10 games. Green = row wins; red = row loses.}
  \label{fig:matrix}
\end{figure}

\paragraph{Reading the table.}
\begin{itemize}
  \item \textbf{All four ``real'' agents crush Random}, as required.
    Random scored 2/40 across the whole tournament -- two draws,
    no wins. This is the baseline check the brief asks for.
  \item \textbf{Stockfish wins 39/40}, confirming it is the right
    reference upper-bound. The single non-loss is a draw vs.\
    Heuristic by threefold repetition.
  \item \textbf{Heuristic, LLM-stub and DRL are roughly comparable
    in the middle.} Heuristic and LLM-stub draw their head-to-head
    1--1 with eight draws -- both play sound positional moves but
    neither has the search depth to consistently break through. DRL
    underperforms Heuristic ($-$2 in 10 games) but beats LLM-stub by
    drawing half its games and winning none -- the value network has
    learned an opening repertoire that is hard to beat but also hard
    to win with on its own.
  \item \textbf{The LLM-stub's hallucinated moves cost it games.}
    Across 40 games the stub had 13 illegal-move emissions caught by
    the parser; 11 were repaired on retry, 2 fell through to the
    heuristic fallback. None lost a game outright (the brief's
    ``invalid move'' clause is never triggered) but they did consume
    moves that could have been productive.
\end{itemize}

\paragraph{Sample game (DRL beats Random as Black, mate in 14).}
\begin{quote}
  \small\ttfamily
  1.~a4 e5 \quad 2.~Na3 Bxa3 \quad 3.~g3 Nf6 \quad 4.~Bg2 O-O \quad
  5.~Bd5 Nxd5 \quad 6.~e3 Bxb2 \quad 7.~Nf3 Bxa1 \quad 8.~Kf1 Bc3 \quad
  9.~Qe2 Qf6 \quad 10.~Ng1 e4 \quad 11.~Qd1 Nc6 \quad 12.~dxc3 Nxc3 \quad
  13.~Ne2 Nxd1 \quad 14.~Ba3 Qxf2\#
\end{quote}

\begin{figure}[h]
  \centering
  \includegraphics[width=0.45\textwidth]{chess_ext/plots/sample_game_final.png}
  \caption{Final position of the sample DRL win: Black's queen on f2
  delivers checkmate; the white king on f1 has no escape squares
  because the black knight on d1 covers e2 and the c-pawn on c2 is
  pinned by the bishop on a3 only nominally -- there is no legal
  defence.}
  \label{fig:samplegame}
\end{figure}

\subsection{Playing online: lichess deployment}
\label{sec:online}

All four of our agents speak the Universal Chess Interface (UCI) via
\texttt{python -m chess\_ext.uci --agent <name>}. To play live games
against humans on lichess.org, the standard recipe is:
\begin{enumerate}
  \item Sign up for a brand-new lichess account, mint an OAuth token
    with the \texttt{bot:play} scope, and convert the account to a BOT.
  \item Clone \texttt{lichess-bot-devs/lichess-bot} \cite{lichessbot}
    and configure \texttt{config.yml} to launch our UCI shim:
\begin{lstlisting}[language=bash]
engine:
  dir: "../lunar_lander_drl"
  name: "uci_runner.sh"
  protocol: "uci"
\end{lstlisting}
    where \texttt{uci\_runner.sh} simply does
\begin{lstlisting}[language=bash]
exec python -m chess_ext.uci \
    --agent drl --model chess_ext/models/drl_value.pt
\end{lstlisting}
  \item Run \texttt{python lichess-bot.py}; humans can now challenge
    your bot.
\end{enumerate}
Full instructions live in \texttt{chess\_ext/PLAYING\_ONLINE.md}.
We did \emph{not} attempt to play \texttt{chess.com} for the reasons
explained at the top of this section.

% =====================================================================
\newpage
\begin{thebibliography}{99}\small

\bibitem{vanhasselt2016}
H.\ van Hasselt, A.\ Guez, D.\ Silver.
\newblock Deep Reinforcement Learning with Double Q-learning.
\newblock \emph{AAAI}, 2016.

\bibitem{wang2016}
Z.\ Wang, T.\ Schaul, M.\ Hessel, et al.
\newblock Dueling Network Architectures for Deep Reinforcement Learning.
\newblock \emph{ICML}, 2016.

\bibitem{hessel2018}
M.\ Hessel, J.\ Modayil, H.\ van Hasselt, et al.
\newblock Rainbow: Combining Improvements in Deep Reinforcement Learning.
\newblock \emph{AAAI}, 2018.

\bibitem{fortunato2018}
M.\ Fortunato, M.\ G.\ Azar, B.\ Piot, et al.
\newblock Noisy Networks for Exploration.
\newblock \emph{ICLR}, 2018.

\bibitem{pathak2017}
D.\ Pathak, P.\ Agrawal, A.\ A.\ Efros, T.\ Darrell.
\newblock Curiosity-driven Exploration by Self-supervised Prediction.
\newblock \emph{ICML}, 2017.

\bibitem{burda2019rnd}
Y.\ Burda, H.\ Edwards, A.\ Storkey, O.\ Klimov.
\newblock Exploration by Random Network Distillation.
\newblock \emph{ICLR}, 2019.

\bibitem{sutton2018}
R.\ S.\ Sutton, A.\ G.\ Barto.
\newblock \emph{Reinforcement Learning: An Introduction (2nd ed.)}.
\newblock MIT Press, 2018.

\bibitem{gymnasium}
Farama Foundation.
\newblock Gymnasium documentation: LunarLander.
\newblock \url{https://gymnasium.farama.org/environments/box2d/lunar_lander/}.

\bibitem{silver2018alphazero}
D.\ Silver, T.\ Hubert, J.\ Schrittwieser, et al.
\newblock A general reinforcement learning algorithm that masters chess,
shogi, and Go through self-play.
\newblock \emph{Science}, 362(6419):1140--1144, 2018.

\bibitem{feng2023chessgpt}
X.\ Feng, Y.\ Luo, Z.\ Wang, et al.
\newblock ChessGPT: Bridging Policy Learning and Language Modeling.
\newblock \emph{NeurIPS}, 2023.

\bibitem{kolasani2025llmchess}
S.\ Kolasani, M.\ Saplin, N.\ Crispino, et al.
\newblock LLM Chess: Benchmarking Reasoning and Instruction-Following
in LLMs through Chess.
\newblock arXiv:2512.01992, 2025.
\newblock LLM Chess leaderboard: \url{https://maxim-saplin.github.io/llm_chess/}.

\bibitem{carlini2023chessllm}
N.\ Carlini.
\newblock Playing chess with large language models.
\newblock \url{https://nicholas.carlini.com/writing/2023/chess-llm.html}, 2023.

\bibitem{strategicrl2025}
\emph{Can Large Language Models Develop Strategic Reasoning?
Post-training Insights from Learning Chess.}
\newblock arXiv:2507.00726, 2025.

\bibitem{zjeffer2022chesszero}
\emph{chess-deep-rl: A research project to create a chess engine
using deep reinforcement learning.}
\newblock \url{https://github.com/zjeffer/chess-deep-rl}.

\bibitem{pythonchess}
N.\ Fiekas.
\newblock python-chess: a chess library for Python.
\newblock \url{https://python-chess.readthedocs.io/}.

\bibitem{lichessbot}
lichess-bot-devs.
\newblock A bridge between Lichess bots and chess engines.
\newblock \url{https://github.com/lichess-bot-devs/lichess-bot}.

\bibitem{stockfish}
The Stockfish developers.
\newblock Stockfish 16: a strong open-source chess engine.
\newblock \url{https://stockfishchess.org/}.

\end{thebibliography}

\end{document}
